\subsection{Applications: II. Functions of measures}

Let \(\mu_1, \mu_2, \ldots, \mu_n\) be real measures on \(T\) , and let \(u(x_1, \ldots, x_n)\) be a finite numerical function, defined on \(\mathbf{R}^n\) , and positively homogeneous, that is (Ch. I, §1, No.~1), such that

\[
u (\alpha x _ {1}, \ldots , \alpha x _ {n}) = \alpha u (x _ {1}, \ldots , x _ {n})\tag{12}
\]

for every scalar \(\alpha \geqslant 0\) . There exist positive measures \(\lambda\) on T such that \(|\mu_i| \leqslant \lambda\) for \(1 \leqslant i \leqslant n\) (for example, the sum \(\sum_{i=1}^{n} |\mu_i|\) ). Let \(\lambda\) and \(\lambda'\) be two such measures on T. One can write \(\mu_i = f_i \cdot \lambda = f_i' \cdot \lambda'\) , where \(f_i\) (resp. \(f_i'\) ) is measurable and essentially bounded for the measure \(\lambda\) (resp. \(\lambda'\) ) (No.~5, Th. 2). We are going to establish the following result: in order that the numerical function \(u(f_1, \ldots, f_n)\) be locally integrable for \(\lambda\) , it is necessary and sufficient that the function \(u(f_1', \ldots, f_n')\) be locally integrable for \(\lambda'\) , in which case

\[
u (f _ {1}, \dots , f _ {n}) \cdot \lambda = u (f _ {1} ^ {\prime}, \dots , f _ {n} ^ {\prime}) \cdot \lambda^ {\prime}.
\]

Since \(|\mu_i| \leqslant \inf(\lambda, \lambda')\) , we can restrict ourselves to the case that \(\lambda \leqslant \lambda'\) . Then \(\lambda = g \cdot \lambda'\) , where \(g\) is a \(\lambda'\) -measurable function such that \(0 \leqslant g \leqslant 1\) (No.~5, Th. 2); whence (No.~4, Prop. 8)

\[
\mu_ {i} = f _ {i} \cdot (g \cdot \lambda^ {\prime}) = (f _ {i} g) \cdot \lambda^ {\prime};
\]

it follows (No.~3, Cor. 2 of Prop. 3) that \(f_{i}g\) is equal to \(f_{i}'\) locally almost everywhere for \(\lambda'\) . Consequently, by (12),

\[
u (f _ {1} ^ {\prime}, \dots , f _ {n} ^ {\prime}) = u (f _ {1} g, \dots , f _ {n} g) = u (f _ {1}, \dots , f _ {n}) g
\]

locally almost everywhere for \(\lambda'\) . In order that \(u(f_1', \ldots, f_n')\) be locally \(\lambda'\) -integrable, it is therefore necessary and sufficient that \(u(f_1, \ldots, f_n)g\) be locally integrable for \(\lambda'\) , therefore (No.~4, Prop. 8) that \(u(f_1, \ldots, f_n)\) be locally integrable for \(\lambda\) ; and then (No.~4, Prop. 8)

\[
u \left(f _ {1} ^ {\prime}, \dots , f _ {n} ^ {\prime}\right) \cdot \lambda^ {\prime} = \left(u \left(f _ {1}, \dots , f _ {n}\right) g\right) \cdot \lambda^ {\prime} = u \left(f _ {1}, \dots , f _ {n}\right) \cdot \lambda .
\]

Thus, the measure \(u(f_{1},\ldots,f_{n})\cdot\lambda\) depends only on the measures \(\mu_{1},\ldots,\mu_{n}\) and the function u; it is also denoted \(u(\mu_{1},\ldots,\mu_{n})\) . This measure is therefore defined whenever u is a positively homogeneous function such that, for a positive measure \(\lambda\) that is an upper bound for all of the \(|\mu_{i}|\) , \(u(f_{1},\ldots,f_{n})\) is locally \(\lambda\) -integrable, where \(f_{i}\) denotes the density of \(\mu_{i}\) with respect to \(\lambda\) . Note that this condition is fulfilled when u is positively homogeneous and continuous: for, one then has

\[
\left| u \left(x _ {1}, \dots , x _ {n}\right) \right| \leqslant a \left(\left| x _ {1} \right| + \left| x _ {2} \right| + \dots + \left| x _ {n} \right|\right)
\]

(u being bounded in a sufficiently small neighborhood of \((0,\ldots,0)\) ), and since \(u(f_{1},\ldots,f_{n})\) is \(\lambda\) -measurable (Ch. IV, §5, No.~3, Th. 1), it is locally \(\lambda\) -integrable by virtue of the criterion for integrability (Ch. IV, §5, No.~6, Th. 5).

Let \(u_1, \ldots, u_p\) be positively homogeneous numerical functions defined on \(\mathbf{R}^n\) , such that the \(p\) functions \(g_k = u_k(f_1, \ldots, f_n)\) ( \(1 \leqslant k \leqslant p\) ) are locally \(\lambda\) -integrable. Let \(v\) be a positively homogeneous numerical function defined on \(\mathbf{R}^p\) such that \(v(g_1, \ldots, g_p)\) is locally \(\lambda\) -integrable. Set

\[
w (x _ {1}, \ldots , x _ {n}) = v \big (u _ {1} (x _ {1}, \ldots , x _ {n}), \ldots , u _ {p} (x _ {1}, \ldots , x _ {n}) \big).
\]

Then, the function w is positively homogeneous, \(w(f_{1}, \ldots, f_{n})\) is locally \(\lambda\) -integrable and, by definition,

\[
w \left(\mu_ {1}, \dots , \mu_ {n}\right) = v \left(u _ {1} \left(\mu_ {1}, \dots , \mu_ {n}\right), \dots , u _ {p} \left(\mu_ {1}, \dots , \mu_ {n}\right)\right).
\]

In the special case of the functions \(x^{+}\) , \(x^{-}\) , \(|x|\) , \(x + y\) , \(\inf(x, y)\) , \(\sup(x, y)\) , the measures defined by the procedure just described coincide, respectively, with those that have been denoted \(\mu^{+}\) , \(\mu^{-}\) , \(|\mu|\) , \(\mu + \nu\) , \(\inf(\mu, \nu)\) , \(\sup(\mu, \nu)\) ; this follows at once from the Cor. of Prop. 2 of No.~2. If \(\mu\) and \(\nu\) are two real measures, and \(\theta = \mu + i\nu\) , then \(|\theta| = \sqrt{\mu^{2} + \nu^{2}}\) ; for, let \(\lambda\) be a measure \(\geqslant 0\) that is an upper bound for \(|\mu|\) and \(|\nu|\) , and let f, g be locally \(\lambda\) -integrable functions such that \(\mu = f \cdot \lambda\) , \(\nu = g \cdot \lambda\) ; then

\[
\sqrt {\mu^ {2} + \nu^ {2}} = \sqrt {f ^ {2} + g ^ {2}} \cdot \lambda ,
\]

\(\theta = (f + ig) \cdot \lambda\) , therefore (No.~2, Prop. 2) \(|\theta| = \sqrt{f^2 + g^2} \cdot \lambda\) .

This method can be applied to the positively homogeneous function \((x_1^2 +\ldots +x_n^2)^{1 / 2}\) to define the length of a curve in \(\mathbf{R}^n\)

